\section{Discussion and Limitations}
\label{sec:discussion}

\paragraph{Relational support.}
Witnesses make the connections supporting a planned answer explicit.
Their validity with respect to the original question nevertheless
depends on the extracted facts and the planner's interpretation.
Semantic matching and source verification accommodate natural-language
variation but remain approximate. Likewise, support for individual
members of an answer set does not establish that the set is complete.
These limitations motivate evaluating both answer quality and the
evidence supporting each answer.

\paragraph{Temporal interpretation.}
Reference-relative ranking allows retrieval to adapt to the period
relevant to a question. Its effectiveness depends on the dates inferred
from the dialogue, particularly when the event time is implicit.
Temporal proximity is a ranking preference rather than a complete
temporal reasoning formalism; ordering and comparison still rely on
the evidence available to the reader.

\paragraph{Evidence and context budget.}
Propositions provide a compact representation of the history, while
passage summaries retain surrounding context. This compression creates
a trade-off between coverage and detail. In particular, selection under
a fact budget can omit part of an accepted witness, so finding a witness
does not guarantee its complete delivery. Evaluation should therefore
consider the actual reader context and token consumption alongside
retrieval and answer metrics. Extraction and reusable summaries also
introduce costs that are shared across questions, making amortized cost
an important part of this comparison.
